\chapter{Canonical Field Definitions}
\label{app:fields}

\section{Six Canonical Critical Point Fields}
\label{sec:appA:canonical}

% Ported from: cwaight_systems_paper.tex
Six vector field environments centered at the origin were used to evaluate the first-order primitives of Chapter~\ref{ch:first_order}. Let $r = \sqrt{x^2 + y^2}$ and $\theta = \text{atan2}(y, x)$. The base fields are the vortex (pure rotational flow), the sink (radial inward flow) whose negation $\mathbf{v}_{\text{source}} = -\mathbf{v}_{\text{sink}}$ is the source (radial outward flow), and the saddle, an axis-aligned saddle rotated by $45^\circ$ so that its stable and unstable directions lie along the diagonals $y = x$ and $y = -x$:
\begin{equation}\label{eq:appA:14}
\mathbf{v}_{\text{vortex}} = \begin{bmatrix} -r \sin(\theta) \\ r \cos(\theta) \end{bmatrix}, \;\;
\mathbf{v}_{\text{sink}} = \begin{bmatrix} -x \\ -y \end{bmatrix}, \;\;
\mathbf{v}_{\text{saddle}} = \begin{bmatrix} -y \\ -x \end{bmatrix}
\end{equation}

% Ported from: cwaight_systems_paper.tex
The sinking vortex (spiraling inward) and spewing vortex (spiraling outward) are linear combinations of the base fields:
\begin{equation}\label{eq:appA:15}
\begin{aligned}
\mathbf{v}_{\text{sink-vortex}} &= 0.4\,\mathbf{v}_{\text{vortex}} + 0.15\,\mathbf{v}_{\text{sink}} \\
\mathbf{v}_{\text{spew-vortex}} &= -0.4\,\mathbf{v}_{\text{vortex}} - 0.15\,\mathbf{v}_{\text{sink}}
\end{aligned}
\end{equation}

\section{Testbed Vortex Fields}
\label{sec:appA:testbed}

% New text
The fixed and sinking vortex maps of Chapter~\ref{ch:zeroth} are given by
(\ref{eq:ch5:rtheta})--(\ref{eq:ch5:sinking}). The fixed vortex there is
the linear clockwise rotation $(y - c_y,\, -(x - c_x))$, the vortex of
(\ref{eq:appA:14}) with the opposite sense. Its sinking component falls
off as $1/r$, so it differs from the linear sinking vortex of
(\ref{eq:appA:15}) and is singular at the center apart from the
$10^{-6}$ regularization.

\section{Steady Double Gyre}
\label{sec:appA:dg}

% New text
The steady double gyre of Chapter~\ref{ch:second_order} is defined in
(\ref{eq:ch3:dg_u}), with its determinant, Hessian, and smaller strain
eigenvalue in closed form in Section~\ref{sec:ch3:dg_example}.

\begin{table}[htbp]
\centering
\caption{Fields used in this dissertation.}
\label{tab:appA:summary}
\begin{tabular}{lll}
\hline
Field & Critical points & Chapter \\
\hline
Vortex, sink, source, saddle & center, stable node, unstable node, saddle & \ref{ch:first_order} \\
Sinking and spewing vortex (linear) & stable and unstable spiral & \ref{ch:first_order} \\
Fixed vortex (testbed) & center & \ref{ch:zeroth} \\
Sinking vortex (testbed, $1/r$ sink) & singular sink at the center & \ref{ch:zeroth} \\
Steady double gyre & two saddles, two centers & \ref{ch:second_order} \\
Santa Barbara Channel record & measured, time-varying & \ref{ch:second_order} \\
\hline
\end{tabular}
\end{table}
